\documentclass[12pt,reqno]{amsart}
\usepackage{math-review}
% Project-wide mathematical notation. No claims are certified here.

\title{The Schur Complement -- Thematic}
\author{Known-mathematics exposition benchmark}
\date{}
\begin{document}
\maketitle
\section{The question and its principal answer}
A coupled quadratic form suggests two related questions: how to test its
positivity, and what remains after eliminating one group of variables.
The same residual matrix answers both, but under the same sign hypothesis.
\begin{theorem}[Classical Schur-complement criterion]\label{thm:main}
Let $p,q\geq1$ be integers. Let $A\in\mathbb R^{p\times p}$ and
$C\in\mathbb R^{q\times q}$ be symmetric, and let $B\in\mathbb R^{p\times q}$.
Assume that $x^{\mathsf T}Ax>0$ for every $0\ne x\in\mathbb R^p$. Define
\[
 M=\begin{pmatrix}A&B\\ B^{\mathsf T}&C\end{pmatrix},
 \qquad S=C-B^{\mathsf T}A^{-1}B.
\]
Then $M$ is positive definite if and only if $S$ is positive definite.
Here a real symmetric matrix $H$ is positive definite when
$v^{\mathsf T}Hv>0$ for every nonzero real vector $v$ of the appropriate size.
\end{theorem}
Throughout the comparison, the unknown test vector is $(x,y)$ and the
blocks $A,B,C$ are fixed. We do not change the field or the meaning of positivity.

\section{The obstruction and square completion}
The term $2x^{\mathsf T}By$ prevents a separate test of the two blocks.
To cancel it, shift only the $x$ variable by a multiple of $y$.
\begin{lemma}\label{lem:square}
Let $A\in\mathbb R^{p\times p}$ be symmetric and invertible,
$B\in\mathbb R^{p\times q}$ and $C\in\mathbb R^{q\times q}$ symmetric.
For all $x\in\mathbb R^p$ and $y\in\mathbb R^q$, putting
$z=x+A^{-1}By$ gives
\[
 x^{\mathsf T}Ax+2x^{\mathsf T}By+y^{\mathsf T}Cy
 =z^{\mathsf T}Az+y^{\mathsf T}(C-B^{\mathsf T}A^{-1}B)y.
\]
\end{lemma}
Denote the left side of this identity by $Q(x,y)$.

This is not yet a positivity proof: the change must preserve nonzero pairs.
\begin{lemma}\label{lem:inverse}
For every real matrix $K\in\mathbb R^{p\times q}$, the linear map
\[
 (x,y)\longmapsto (x+Ky,y)
\]
is invertible, with inverse $(z,y)\mapsto(z-Ky,y)$. Consequently it takes
nonzero vectors to nonzero vectors, in both directions.
\end{lemma}
With $K=A^{-1}B$, the two lemmas prove sufficiency because
$z^{\mathsf T}Az+y^{\mathsf T}Sy>0$ for every $(z,y)\ne0$ when $A,S>0$.
For necessity, choose $x=-A^{-1}By$ for each nonzero $y$. Then
\[
 y^{\mathsf T}Sy=Q(-A^{-1}By,y)>0.
\]
The same calculation has a variational interpretation. For fixed $y$,
strict positivity of $A$ gives a unique minimum:
\[
 \min_{x\in\mathbb R^p}Q(x,y)=y^{\mathsf T}Sy,
 \qquad x_{\min}=-A^{-1}By.
\]
Thus the Schur complement is the residual quadratic form after eliminating $x$.

\section{Congruence: the same elimination in coordinates}
The second description works directly with block matrices rather than
with the expanded quadratic form.
\begin{proposition}\label{prop:comparison}
Let $A\in\mathbb R^{p\times p}$ be symmetric and invertible,
$B\in\mathbb R^{p\times q}$ and $C\in\mathbb R^{q\times q}$ symmetric.
For
\[
 M=\begin{pmatrix}A&B\\B^{\mathsf T}&C\end{pmatrix},\quad
 T=\begin{pmatrix}I_p&-A^{-1}B\\0&I_q\end{pmatrix},\quad
 S=C-B^{\mathsf T}A^{-1}B,
\]
one has $T^{\mathsf T}MT=\operatorname{diag}(A,S)$.
Because $T$ is invertible, $M>0$ is equivalent to $A>0$ and $S>0$ jointly;
$S>0$ alone is not sufficient without $A>0$.
\end{proposition}
The calculation can be checked without a new existence theorem:
\[
 MT=\begin{pmatrix}A&0\\B^{\mathsf T}&S\end{pmatrix},\qquad
 T^{\mathsf T}MT=\begin{pmatrix}A&0\\0&S\end{pmatrix}.
\]
For $v=(z,y)$, the vector $Tv=(z-A^{-1}By,y)$ is exactly the inverse
change used above. Hence the congruence identity is square completion
written before evaluating on a vector, not a stronger criterion.

\section{Comparison, scope and a failure boundary}
The viewpoints share the same data and conclusions. Square completion
explains the choice of shift; congruence packages that shift as an invertible
matrix; minimization identifies the residual value when $A>0$.

Their usable interfaces differ. The congruence and square identities need
only symmetric invertible $A$, whereas the displayed global minimum requires
$A>0$. Thus the identities survive a loss of positivity, but the sign conclusion does not.
\[
 A=-1,\quad B=0,\quad C=1
 \quad\Longrightarrow\quad S=1,
 \qquad M=\operatorname{diag}(-1,1).
\]
This is a counterexample to the stronger claim that $S>0$ alone suffices,
not to the stated theorem. Singular $A$ is outside all formulas using $A^{-1}$;
no pseudoinverse extension or historical novelty claim is intended.

\nocite{S-guide}
\bibliographystyle{plain}
\bibliography{../generated/references}
\end{document}
